\section{与原文的问答录：六十个有原文依据的问题}\label{app:dialogue}

本附录把原文里\textbf{明确写明}的内容整理成对话体：读者会问的问题、依据原文的答案、
以及原文的短引。每条都带印刷页锚点，可以逐条回原文核对。它与附录 B 的推导链语料互补：
附录 B 记的是“我们怎么把这一步做出来”，本附录记的是“原文到底怎么说的”。

选取标准只有一条：\textbf{只收原文明确陈述的内容}（A 级）。原文只是暗示、或属于我们的推断，
一律不收——那些放在正文里并标 \pB / \pC。因此本附录可以当作一份“最小事实清单”：
如果读者只想先知道原文说了什么，读完这三节就够了。

用法：读正文卡住时，先在这里查有没有对应的问题；讨论班可以把本附录当“原文核对”环节的题卡——
先答，再回原文对页码。

% ======================================================================
%  53a_dialogue_A.tex — 附录 D 第一部分《舞台与机器》问答录
%  底本：KW 2006（hep-th/0604151），§1--§5，印刷 pp.2--77
%  只收原文明确陈述的 A 级内容；每条答案可回原文核对，短引 <=120 字符
% ======================================================================

\subsection{舞台与机器：\S1--\S5（pp.2--77）}

以下二十问按原文顺序走一遍 \S1--\S5：从 $\mathcal N=4$ 超对称 Yang--Mills 的场内容与 $S$-对偶，经过扭化、以 $\PP^1$ 参数化的一族拓扑场论与典范参数 $\Psi$，到紧化给出的 Hitchin 模空间 $\MH$，最后到二维拓扑场论、generalized complex geometry，以及 $S$-对偶的镜像对称解释。答案只复述原文明确写下的内容，并附原文短引与印刷页锚点，可逐条回原文核对。原文只是暗示、没有明说的，不在此列。

\QAq{这篇论文的总目标是什么？}
\QAa{论文的目标是用四维超对称规范理论来理解复 Riemann 面的几何 Langlands 对应；原文在导言中把这套计划概括为六个主要想法，其中一个想法正是：在参数的某个特殊值上，四维 $S$-对偶表现为二维镜像对称。}
\QAs{src p.2}{``Our focus in the present paper is on the geometric Langlands program for complex Riemann surfaces.''}

\QAq{$\mathcal N=4$ 超对称 Yang--Mills 是怎么构造的？场有哪些？}
\QAa{原文从十维超对称 Yang--Mills 作维数约化构造它：十维的规范场 $A$（$G$-丛 $E$ 上的联络）与费米场 $\lambda$ 约化到四维后，$A_\mu$ 仍是规范场，而 $A_I$（$I\ge 4$）成为六个标量 $\varphi_i=A_{i+4}$。原文随后把这六个标量与一个 $G$-丛上的联络作为四维理论的场内容。}
\QAs{src p.6}{``N = 4 super Yang-Mills is most easily constructed by dimensional reduction from ten dimensions.''}

\QAq{$S$-对偶的核心断言是什么？}
\QAa{原文断言存在一个额外的量子对称性：它反演复耦合 $\tau$，把 $G$ 换成 Langlands 对偶群 $\Ldual$，并交换电与磁的电荷。对 simply-laced 的 $G$，$\tau\to-1/\tau$ 与经典的 $\tau\to\tau+1$ 生成 $SL(2,\ZZ)$；非 simply-laced 时第二生成元改为 $\tau\to-1/n_g\tau$。}
\QAs{src p.11}{``an additional quantum symmetry that inverts $\tau$, exchanges $G$ with $\Ldual$''}

\QAq{“扭化”在做什么？为什么选 GL 扭化？}
\QAa{扭化是指用 $\mathrm{Spin}(4)\times\mathrm{Spin}(6)$ 的另一个子群 $\mathrm{Spin}'(4)$ 替换原来的 $\mathrm{Spin}(4)$，使自旋表示中出现不变向量，从而留下一个不变超荷。原文指出，在 $\mathcal N=4$ 理论可用的几种扭化里，与几何 Langlands 对应相关的那一种被称为 GL 扭化。}
\QAs{src p.17}{``It turns out to be the twist relevant to the geometric Langlands program, and we will call it the GL twist.''}

\QAq{不变超荷 $Q$ 从哪里来？为什么 $Q^2=0$？}
\QAa{在 $S_+$ 中任取一个 $\mathrm{Spin}'(4)$ 不变向量，就得到一个不变超荷 $Q$；由于超对称代数满足 $\{Q,Q\}=\{\bar Q,\bar Q\}=0$，它自动满足 $Q^2=0$。原文由此把物理问题换成对 $Q$-上同调的研究。}
\QAs{src p.18}{``It will automatically obey $Q^2 = 0$, by virtue of (2.15).''}

\QAq{扭化后玻色场变成了什么？}
\QAa{原文说扭化后的玻色场是一个规范场，局域地写为取值在伴随表示中的 one-form $A=A_\mu dx^\mu$，再加上第二个伴随值 one-form $\varphi=\varphi_\mu dx^\mu$，以及复标量 $\sigma,\bar\sigma$。原来的六个标量中，$\varphi_0,\dots,\varphi_3$ 组成这个 one-form，$\sigma,\bar\sigma$ 由 $\varphi_4,\varphi_5$ 组合而成。}
\QAs{src p.19}{``a gauge field, which locally is an adjoint-valued one-form ... along with a second adjoint-valued one-form''}

\QAq{为什么会出现一族以 $\PP^1$ 参数化的拓扑场论？}
\QAa{不变超荷取 $\epsilon=u\epsilon_\ell+v\epsilon_r$，只有比值 $t=v/u$ 有意义，于是所有选择构成一个 $\PP^1$；原文说这就给出以 $\PP^1$ 参数化的一族拓扑场论。$t$ 有自然的仿射含义，且 $t$ 与 $-t$ 给出等价的拓扑场论，后者来自 $\mathrm{SU}(4)_R$ 中心的作用。}
\QAs{src p.22}{``Thus, the topological field theories with parameter $t$ or $-t$ are equivalent.''}

\QAq{这些拓扑场论的超对称方程是什么？$t=\pm i$ 时为什么变成平坦联络？}
\QAa{要求所有费米场变分为零，得到 $(F-\varphi\wedge\varphi+tD\varphi)^+=0$、$(F-\varphi\wedge\varphi-t^{-1}D\varphi)^-=0$ 与 $D*\varphi=0$。在 $t=\pm i$ 时令 $\mathcal A=A+i\varphi$，前两式等价于 $\mathcal F=0$，即 $\mathcal A$ 是平坦的 $G_{\mathbb C}$ 联络，从而给出同态 $\pi_1(M)\to G_{\mathbb C}$。}
\QAs{src p.25}{``a solution of (3.29) at $t = \pm i$ determines a complex-valued flat connection''}

\QAq{为什么说这个理论是拓扑场论？}
\QAa{原文的论证是：作用量 $I_0$ 形如 $\delta_\ell\delta_r V$，而度规不进入 $\delta_\ell,\delta_r$ 的定义，只进入 $V$ 的选取。因此作用量对度规的依赖是 $\delta_\ell\delta_r(\cdots)$ 的形式，这正是导向拓扑场论的关键性质。}
\QAs{src p.33}{``it also has the key property that leads to a topological field theory''}

\QAq{典范参数 $\Psi$ 是什么？为什么理论只依赖它？}
\QAa{原文说理论只依赖 $t$ 与 $\tau$ 的一个组合，称之为典范参数：$\Psi=\frac{\tau+\bar\tau}{2}+\frac{\tau-\bar\tau}{2}\cdot\frac{t-t^{-1}}{t+t^{-1}}$。把作用量写成 $I=\delta_tV+\frac{i\Psi}{4\pi}\int\Tr F\wedge F$ 后，$\Psi$ 以外的依赖都被 $\delta_t$ 吸收，所以只有 $\Psi$ 是重要的。}
\QAs{src p.35}{``only depends on a certain combination $\Psi$ of $t$ and $\tau$, which we will call the canonical parameter''}

\QAq{$\Psi$ 有哪些特殊值？}
\QAa{原文列出：$t=\pm 1$ 时 $\Psi=\Re\tau$；一般地 $|t|=1$ 时 $\Psi$ 为实数；$t=\pm i$ 时 $\Psi=\infty$，与 $\tau$ 无关。反过来 $t^2=-(\Psi-\bar\tau)/(\Psi-\tau)$，故任意 $\Psi$ 都能在任意弱的耦合下实现。}
\QAs{src p.36}{``If $t = \pm i$, then $\Psi = \infty$, independent of $\tau$.''}

\QAq{把四维理论紧化到 $\Sigma\times C$ 上会发生什么？}
\QAa{取 $C$ 为亏格大于 $1$ 的紧 Riemann 面、$\Sigma$ 远大于 $C$，则扭化的四维 $\mathcal N=4$ 规范理论在 $\Sigma$ 上约化为一个 sigma 模型，其映射为 $\Phi:\Sigma\to\MH$，其中 $\MH$ 是 $C$ 上 Hitchin 方程的解模空间。原文随后指出，长距离的有效物理就是目标空间为 $\MH(G,C)$ 的超对称 sigma 模型。}
\QAs{src p.38}{``reduces on $\Sigma$ to a sigma-model of maps $\Phi : \Sigma \to \MH$''}

\QAq{为什么 $\MH$ 是超凯勒流形？}
\QAa{原文把 Hitchin 方程等同于三个矩映射的消失，从而 $\MH$ 是无穷维仿射空间 $W$（所有 $(A,\varphi)$）在规范变换群下的超凯勒商。超凯勒商保持超凯勒结构，所以 $\MH$ 是超凯勒流形；原文也用这一点解释二维一侧的 $(4,4)$ 超对称。}
\QAs{src p.46}{``$\MH$ is the hyper-Kahler quotient of the infinite-dimensional space $W$ by the group of gauge transformations''}

\QAq{$\MH$ 的三个复结构 $I,J,K$ 分别看到什么？}
\QAa{$\MH$ 上有满足 $IJ=-JI=K$ 及循环置换的三个复结构 $I,J,K$，还有以 $\PP^1_h$ 参数化的一族复结构 $I_w$，其中 $J=I_{-i}$、$K=I_1$。在复结构 $I$ 下全纯数据是 Hitchin 对 $(E,\varphi)$；在复结构 $J$ 下，$\MH$ 是稳定同态 $\vartheta:\pi_1(C)\to G_{\mathbb C}$ 的模空间 $Y$（Corlette--Donaldson 定理）。}
\QAs{src p.51}{``in complex structure $J$, is the moduli space $Y$ that parametrizes stable homomorphisms''}

\QAq{Hitchin 纤维化是什么？底空间与纤维有多大？}
\QAa{把 Hitchin 对 $(E,\varphi)$ 送到 $w=\Tr\varphi^2$（一般 $G$ 改用独立的 Casimir 不变量），得到映到复向量空间 $B$ 的 Hitchin 纤维化 $\pi:\MH\to B$。$B$ 的复维是 $(g-1)\dim G$，等于 $\MH$ 复维的一半；纤维 $F$ 的维数与 $B$ 相同，且 Hitchin 纤维化是满射。}
\QAs{src p.53}{``The Hitchin fibration, as it is most commonly called, is the map $\pi : \MH \to B$''}

\QAq{二维拓扑场论什么时候是 A 模型、什么时候是 B 模型？}
\QAa{二维拓扑场论由一对复结构 $(J_+,J_-)$ 决定：$J_+=J_-$ 时是 B 模型，$J_+=-J_-$ 时是 A 模型，其余情形可用 generalized complex geometry 化到 A 模型。从四维约化看 $(w_+,w_-)=(-t,t^{-1})$，所以 $t=\pm i$ 恰给出 B 模型，$t$ 为实数时给出 A 模型。}
\QAs{src p.61}{``For a B-model, we want $w_+ = w_-$. So the condition is $t^{-1} = -t$, which occurs precisely for $t = \pm i$.''}

\QAq{generalized complex geometry 在这里起什么作用？}
\QAa{原文用它把二维拓扑场论写成广义复结构：由 $\MH$ 的超凯勒度量与一对点 $J_+,J_-\in\PP^1_h$ 得到一对广义复结构，其中决定模型的那个只依赖 $J$、$\omega_K$ 与 $\omega_I$，这三个量与 $C$ 的复结构无关，故约化得到的这族理论不依赖 $C$ 的度规。进一步，一般情形下它是 A 模型再加一个 $B$ 场变换，由此给出 $\Psi$ 的二维解释。}
\QAs{src p.67}{``The first important observation is that $I$, $K$, and $\omega_J$ have disappeared.''}

\QAq{$S$ 交换 $\Psi=0$ 与 $\Psi=\infty$；这两个端点各是什么模型？}
\QAa{原文取 $\tau$ 为纯虚数：$\Psi=\infty$（$t=i$）给出复结构 $J$ 的 B 模型，与 $\tau$ 的取值无关；$S$ 把 $t$ 送到 $1$，得到复结构 $K$ 的 A 模型，即 $\Psi=0$。因此 $S$ 把 $\MH(G,C)$ 的 B 模型换成 $\MH(\Ldual,C)$ 的 A 模型。}
\QAs{src p.70}{``that maps $\Psi \to -1/n_g \Psi$ and hence exchanges $\Psi = 0$ and $\Psi = \infty$''}

\QAq{$S$-对偶怎样在 Hitchin 纤维化上实现镜像对称？}
\QAa{原文要证的是 $S$-对偶在 Hitchin 纤维化 $\pi:\MH\to B$ 的底空间上经典地作用，而在纤维上作 $T$-对偶。零膜 $B_p$ 是 (B,B,B) 型，$S$ 把它变成支撑在某个纤维 $\tilde F$ 上的 (B,A,A) 型膜，两侧模空间分别是 $F$ 与 $\tilde F$ 上秩一平坦 Chan--Paton 丛的复环面，$S$ 在它们之间建立同构。原文总结说这正是 SYZ 的环面纤维对偶，只是底空间之间的映射 $\Xi:B\to{}^LB$ 在 SYZ 中通常被当作平凡。}
\QAs{src p.71}{``$S$-duality acts classically on the base of the Hitchin fibration''}

\QAq{$\MH$ 上有哪些类型的膜？$S$-对偶怎样交换它们？}
\QAa{原文列出 (B,B,B) 型膜，即在 $I,J,K$ 三个复结构下都是 B 型，以及 (B,A,A)、(A,B,A)、(A,A,B) 型膜，即在一个复结构下是 B 型、在另外两个下是 A 型。$S$ 把 (B,B,B) 型换成 (B,A,A) 型，而 (A,B,A) 与 (A,A,B) 型各自换成同一类型的膜。}
\QAs{src p.77}{``(B, B, B)-branes, which are B-branes in each complex structure $I$, $J$, $K$''}

% ======================================================================
%  53b_dialogue_B.tex — 附录 D 第二部分《线算子与特征膜》问答录
%  底本：KW 2006（hep-th/0604151），§6--§8，印刷 pp.78--122
%  只收原文明确陈述的 A 级内容；每条答案可回原文核对，短引 <=120 字符
% ======================================================================
\subsection{线算子与特征膜：\S6--\S8（pp.78--122）}

本录只收原文 \S6--\S8 明写的内容：线算子与圈算子的定义及其物理意义、紧化到二维后它们变成什么、\S7 的离散通量与 $S$-对偶、以及 \S8 中 Wilson 算子如何作用于膜。问题按原文页序排列，答案紧贴原文措辞；凡属我们的推测、或原文只是暗示而非明说的，一律不收。每条末尾的英文引文逐字取自原文（必要时截成语段），方括号中标出它在印刷本中的页码，可据以回查底本。

\QAq{\S6 为什么把算子挂到一维流形上？}
\QAa{原文说这套四维拓扑场论的特殊之处，在于它把算子关联到 $M$ 中定向的一维流形 $S$ 上，而算子只依赖 $S$ 的同伦类、不只是同调类。作者随即点明动机：重点并不是研究 $M$ 的基本群，而是这些一维流形算子与膜之间表现出很有意思的关系。}
\QAs{src p.78}{``these one-manifold operators prove to have very interesting properties in relation to branes''}

\QAq{什么情形下叫圈算子，什么情形下才有线算子？}
\QAa{原文规定：$M$ 紧致且无边界时，$S$ 必须是闭圈，这类一维流形算子因此通称圈算子。若四维流形非紧，则允许延伸到无穷远的线算子；其中最要紧的是在某个空间点贯穿全部时间的静态线算子，它在量子化时必须被包括进来。}
\QAs{src p.79}{``If $M$ is compact and without boundary, then $S$ is required to be a closed loop''}

\QAq{Wilson 圈算子的定义式是什么？}
\QAa{原文取 $G$ 的不可约表示 $R$ 与 $M$ 中定向的圈 $S$，定义 $W_0(R,S)=\Tr_R P\exp(-\oint_S A)$，即 $A$ 绕 $S$ 的和乐在 $R$ 中的迹。负号来自该书的约定 $D=\dd+A$；若把规范场写成 hermitian 的 $A'=\ii A$，指数就变成 $+\ii\oint A'$。要在扭化理论里保持拓扑对称，须取 $t=\pm\ii$，把 $A$ 换成 $A+\ii\varphi$ 或 $A-\ii\varphi$：此时拓扑方程恰好断言它平坦，其和乐才是自然的不变量。}
\QAs{src p.80}{``$W_0(R,S)$ is the trace, in the representation $R$, of the holonomy of $A$ around $S$''}

\QAq{开曲线上的 Wilson 算子为什么不取迹？}
\QAa{原文给出开线情形的定义：若 $S$ 有端点 $p$、$q$，就把 $W(R,S)$ 定义为平行输运矩阵而非迹，它是 $p$ 处纤维到 $q$ 处纤维的映射；迹只在闭圈情形出现。作者特别指出，这个推广版本在真正进入几何 Langlands 时是不可缺的。}
\QAs{src p.80}{``we define $W(R,S)$ not as a trace but as the matrix of parallel transport''}

\QAq{'t Hooft 算子与 Wilson 算子是什么关系，它怎样定义？}
\QAa{原文把 Wilson 算子称为 order 算子，并说 order 算子的对偶通常是 disorder 算子；在本问题里，Wilson 算子的对偶正是 't~Hooft 算子。在 abelian 的基本情形里，'t~Hooft 算子由路径积分定义：只对 $M\setminus S$ 上方线丛的联络积分，并要求它们在 $S$ 附近与奇异的 Dirac 单极子相符，使曲率在 $S$ 附近按 $F\sim(\ii/2)\star_3\dd(1/r)$ 发散——这就是“用奇异规范场定义”的含义。与 Wilson 需要 $S$ 的定向不同，'t~Hooft 圈算子需要的是 $S$ 法丛的定向；荷 $m$ 的算子就把基本奇点取 $m$ 倍。}
\QAs{src p.82}{``the path integral is to be performed over connections on a line bundle over $R^4\setminus L$''}

\QAq{为什么 't Hooft 奇点必须让 $\varphi$ 也奇异？}
\QAa{因为扭化理论的拓扑方程 $(F+\dd\varphi)^+=(F-\dd\varphi)^-=0$ 在 $F$ 有 Dirac 单极奇点而 $\varphi$ 光滑时无解。原文因此要求场在 $S$ 附近取 (6.9) 的奇点：$F=(\ii/2)\star_3\dd(1/|x|)$、$\varphi=(\ii/2|x|)\dd s$，其中 $s$ 是沿 $S$ 的“时间”坐标；荷 $m$ 即取 $m$ 倍。}
\QAs{src p.83}{``These equations are not satisfied if $F$ has the Dirac monopole singularity and $\varphi$ is non-singular''}

\QAq{任意紧群 $G$ 的 't Hooft 算子怎样分类？}
\QAa{原文任取一个同态 $\rho:U(1)\to G$，要求场在 $S$ 附近的奇点是 (6.9) 在 $\rho$ 下的像；算子只依赖 $\rho$ 的共轭类。共轭类同态由对偶群 $\Ldual$ 的最高权分类，所以每个 $\Ldual$ 的不可约表示 $\LR$ 给出一个 't~Hooft 算子 $T(\LR,S)$，这就是经典的 GNO 对偶。同一条对应还使 't~Hooft 算子的分类等同于 $\PP^1$ 上 $G$-丛的分类。}
\QAs{src p.84}{``The fact that 't~Hooft operators of $G$ are classified by representations of $\Ldual$ is the classic GNO duality''}

\QAq{紧化到二维后，四维的圈算子变成什么？}
\QAa{原文取 $M=\Sigma\times C$（$\Sigma$ 远大于 $C$，先设 $\Sigma$ 无边）。若圈形如 $S=p\times S'$（$p\in\Sigma$、$S'\subset C$），它在 $\Sigma$ 上的有效理论里就是一个点状局域算子；若 $S=S''\times q$，它仍是二维的圈算子；一般的曲线同伦于在 $\Sigma$ 与 $C$ 之间交替传播的曲线，投影到 $\Sigma$ 上便成为带局域算子插入的圈算子。}
\QAs{src p.90}{``simply looks like a pointlike local operator in the effective theory on $\Sigma$''}

\QAq{二维里线算子为什么可以用膜来理解？}
\QAa{原文指出，二维流形被一维流形局部地分成两个不相交的区域，线算子两侧的耦合可以不同，因此它可能产生长程效应。用折叠技巧把两侧放到同一侧后，沿算子折叠出的边界条件就是一个 $(Y,Z)$-膜；而这里所说的“膜”，一般就是指量子场论里的局域边界条件。于是线算子在二维里归结为一个膜。}
\QAs{src p.91}{``the line operator reduces in two dimensions to a $(Y,Z)$-brane''}

\QAq{线算子靠近边界时起什么作用？}
\QAa{原文把线算子 $X$ 与膜 $B$ 合成的有效边界条件记作 $B'=XB$，即 $X$ 作用在旧膜上产生一个新膜。因此可以把线算子看成作用在边界条件上、产生新边界条件的对象。这正是几何 Langlands 关注的情形：$\Sigma$ 有边界，线算子与边界平行。}
\QAs{src p.94}{``we write $B'=XB$ for the new brane that is produced by acting with $X$ on the old brane''}

\QAq{什么叫线算子的特征膜（eigenbrane）？}
\QAa{原文的定义是：若对某个向量空间 $V$ 有 $XB=B\otimes V$，就称 $B$ 是线算子 $X$ 的特征膜。做法是：对每个膜 $B$ 与向量空间 $V$ 都定义一个新膜 $B\otimes V$；本征值应看成一个向量空间而不是数（$V=\CC^k$ 时只是取 $k$ 份）。原文说这个概念正对应 Beilinson--Drinfeld 在几何 Langlands 中定义的特征层。}
\QAs{src p.95}{``A brane $B$ is an eigenbrane for a line operator $X$ if $XB=B\otimes V$, for some vector space $V$''}

\QAq{电特征膜、磁特征膜是什么？$S$-对偶怎样互换二者？}
\QAa{原文定义：Wilson 线算子的联合特征膜叫电特征膜，'t~Hooft 算子的联合特征膜叫磁特征膜。因为四维来的线算子可在 $C$ 上任取一点 $p$，它们自动对易；而 $V_p$ 局部不依赖 $p$、只有整体单值，故拼成 $C$ 上的平坦向量丛——所以电、磁特征膜都自动带着 $C$ 上的平坦向量丛。$S$-变换 $S:\Psi\mapsto-1/(n_g\Psi)$ 把拓扑 Wilson 算子送到拓扑 't~Hooft 算子，因而把电特征膜送到磁特征膜。}
\QAs{src p.97}{``A joint eigenbrane of the Wilson line operators will be called an electric eigenbrane''}

\QAq{离散电通量与离散磁通量是怎样定义的？}
\QAa{原文先指出中心 $Z=\pi_1(G_{\mathrm{ad}})$ 自然自对偶，存在 perfect pairing $\Upsilon:Z\times Z\to U(1)$。在 $M=S^1\times W$ 上，丛的拓扑由 $a\in H^1(W,Z)$ 与 $m\in H^2(W,Z)$ 给出；character $e:H^1(W,Z)\to U(1)$ 与 $m$ 就分别称为离散电通量与离散磁通量，而 $\sum_a e(a)Z_{a,m}$ 可在 Hilbert 空间 $\mathcal H_{e,m}$ 里读成迹。$S$-对偶交换 $E$ 与 $M$，把 $(e,m)$ 送到 $(-m,e)$。}
\QAs{src p.105}{``$e$ and $m$ are called respectively the discrete electric and magnetic flux''}

\QAq{什么叫 universal bundle？它一定存在吗？}
\QAa{原文定义：universal bundle 是 $G$-丛 $\mathcal E\to M\times C$，使每个 $p\in M$ 上 $\mathcal E|_{p\times C}$ 都同构于 $E_p$。对伴随群 $G_{\mathrm{ad}}$ 粘合没有障碍，故 $\mathcal E_{\mathrm{ad}}$ 存在；但对一般 $G$，$\Theta_{ki}\Theta_{jk}\Theta_{ij}=f_{ijk}\in Z$ 给出非平凡元素 $\zeta$，它是 universal $G$-丛存在的阻碍。原文补充：它总可以作为 twisted vector bundle 存在，twist 为 $\varrho(\zeta)$。}
\QAs{src p.109}{``it does always exist as a twisted vector bundle, twisted by $\varrho(\zeta)$''}

\QAq{\S7.2 里 $m_0$ 与 $e_0$ 各是什么？}
\QAa{原文说靶空间 $\MH(G_{\mathrm{ad}},C)$ 并不连通，其连通分支由 $G_{\mathrm{ad}}$-丛 $E\to C$ 的拓扑型标记，$m_0$ 就是这个拓扑型。$e_0$ 则给出 $\MH(G_{\mathrm{ad}},C)$ 上的平坦 $B$-场 $b_{e_0}=e_0(\zeta)$，它在有效 sigma 模型里的作用正是按平坦 $B$-场的样式给每个映射加权；$e_0\ne0$ 时，膜的 Chan-Paton 丛是 twisted vector bundle。}
\QAs{src p.111}{``Its components are labeled by the topological type of the $G_{\mathrm{ad}}$-bundle $E\to C$''}

\QAq{\S7.2 后半的 $e_1$、$m_1$ 度量什么？}
\QAa{原文把 $\Sigma=S^1\times\tilde S^1$ 时 $\xi(E)$ 的其余分量记作 $a_1$、$m_1$：$m_1$ 度量 $\Phi|_{\tilde S^1}$ 在基本群中的同伦类，$a_1$ 度量沿时间方向的绕数。在二维图像里，$\MH(G)$ 上的弦有守恒动量 $E_C=H^1(C,\ZZ)$，$\MH(G_{\mathrm{ad}})$ 上的弦有守恒绕数 $M_C=\pi_1(\MH(G_{\mathrm{ad}}))=H^1(C,\ZZ)$。$S$-对偶交换这两套守恒量，原文说这正是它与二维 $T$-对偶有关的体现。}
\QAs{src p.115}{``S-duality or Montonen-Olive or Langlands duality exchanges the two pictures''}

\QAq{\S8.1 里 Wilson 线怎样作用于膜？}
\QAa{原文结论：表示 $R$、支在点 $p\in C$ 的 Wilson 线把膜 $B$ 的 Chan-Paton 丛 $U$ 变成 $U\otimes E_p(R)$。$E_p(R)$ 一般不是向量丛，而是 twist 为 $\theta_R(\zeta)$ 的 twisted vector bundle，所以这个作用把 $B$-场改成 $b\to b+\theta_R(\zeta)$，也就是把 \S7 的离散电场改成 $e_0\to e_0+\theta_R$。$S$-对偶的像则是 't~Hooft 线把 $m_0$ 改成 $m_0+\xi(\LR)$。}
\QAs{src p.119}{``it changes the discrete electric field studied in section 7 by $e_0\to e_0+\theta_R$''}

\QAq{为什么支在一点的零膜是电特征膜？它与 $\pi_1(C)$ 的表示有什么关系？}
\QAa{原文论证：任何向量丛限制到一点都平凡，故 $E_q(R)$ 限制到 $x\in\MH$ 也平凡，于是 $U_x\otimes E_q(R)=U_x\otimes E_q(R)|_x$；零膜 $B_x$ 就是电特征膜，本征值是向量空间 $E_q(R)|_x$。在几何 Langlands 一侧，起手是 $\Ldual$ 与同态 $\vartheta:\pi_1(C)\to\Ldual_\CC$，它定义 $\MH(\Ldual,C)$ 中的点 $x(\vartheta)$；若 $\vartheta$ 不可约，则 $x(\vartheta)$ 是光滑点，支在它上面的零膜就是电特征膜。原文还指出，电特征膜其实是不同 $e_0$ 的 twisted 零膜之和。}
\QAs{src p.121}{``of the Wilson loop operator $W_q(R)$ being the vector space $E_q(R)|_x$''}

\subsection{Hecke 算子与翻译器：\S9--\S12（pp.123--222）}

本部分覆盖原文第 9 至第 12 节（印刷页 pp.123--222），主题是 't~Hooft 算子、Hecke 修改以及相应 A-膜在四维规范理论中的构造。问答只取原文用陈述句明确给出定义、结论或公式的地方（A 级材料）；凡原文只是暗示、没有明说的推论一律不录。每条答案之后给出印刷页锚点与英文短引（不超过 120 字符），以便逐条回查原文。

\QAq{\S9 为什么值得单独展开？}
\QAa{原文在本节要证明，$\Psi=0$ 处的超对称 't~Hooft 算子恰好对应几何 Langlands 的 Hecke 算子；$\Psi=0$ 由取 $t=1$、$\theta=0$ 达到。此时四维 TQFT 在二维化后成为以 $\MH$ 为目标、复结构 $K$ 的 A-模型；A-模型没有瞬子修正，超对称位形都时间无关。}
\QAs{src p.123}{``'t~Hooft operators correspond to the Hecke operators of the geometric Langlands program.''}

\QAq{'t~Hooft 算子是普通意义下的算子吗？}
\QAa{不是。原文强调静态 't~Hooft 算子支撑在 $\RR\times p$ 上，属于“线算子”，它们必须被并入物理的、超对称的态空间的定义之中，而不能当作通常 Hilbert 空间上的算子。这一定位正是后文把 't~Hooft 算子与 Hecke 修改对应起来的前提。}
\QAs{src p.123}{``they must be incorporated in the definition of the space of physical, supersymmetric states.''}

\QAq{$U(1)$ 情形 't~Hooft 算子怎样作用？}
\QAa{算子 $T(m;p_0)$ 把线丛 $L$ 变为 $L\otimes\mathcal O(p_0)^m$，即按 $\mathcal O(p_0)$ 的 $m$ 次幂作扭化。原文指出这与 $U(1)$ 的 Hecke 算子的标准定义一致；同一 $p_0$ 处的算子构成交换群 $T(m)=T(1)^m$，不同点的算子也彼此交换。}
\QAs{src p.131}{``This result agrees with the standard definition of the Hecke operators for U(1).''}

\QAq{$U(N)$ 情形的奇点与 Hecke 修改如何对应？}
\QAa{'t~Hooft 算子由同态 $\rho:U(1)\to U(N)$（等价于整数权 $m_1\ge\cdots\ge m_N$）的共轭类分类；在 $p_0$ 附近 $E$ 分解为 $L_1\oplus\cdots\oplus L_N$，修改把 $L_i$ 换成 $L_i\otimes\mathcal O(p_0)^{m_i}$。原文说这与 $U(N)$ 的 Hecke 算子的数学描述完全一致，并称之为 't~Hooft 算子与 Hecke 算子之间的基本联系。}
\QAs{src p.132}{``This coincides with the usual mathematical description of the action of the Hecke operators for U(N).''}

\QAq{原文所说的 Hecke 修改究竟指什么？}
\QAa{一个 Hecke 修改是带指定同构的丛对：$E_+$ 在 $p_0$ 之外与 $E_-$ 自然同构，而这一同构一般不能延拓到 $p_0$。局部地看 $E_+=\oplus_i\,(L_i\otimes\mathcal O(p_0)^{m_i})$；不同的局部分解可能给出同一个 $E_+$，原文正是要研究这些分解商掉等价关系之后的空间。}
\QAs{src p.136}{``any Hecke modification $E_+$ of $E_-$ is naturally isomorphic to $E_-$ away from $z = 0$''}

\QAq{Hecke 修改的空间为何需要紧化？}
\QAa{A-模型的物理态是这些修改空间上的上同调，原文指出紧性使相应上同调群的定义简单直接。$U(2)$ 的权 $(1,0)$ 给出光滑紧的 $\CC\PP^1$；而 $(2,0)$ 的自然紧化必须额外允许 $(1,1)$ 型修改，后者就是 $E_-\otimes\mathcal O(p_0)$。}
\QAs{src p.138}{``Compactness makes the definition of the appropriate cohomology groups much more straightforward''}

\QAq{仿射 Grassmannian 与 Hecke 修改空间是什么关系？}
\QAa{带 $p_0$ 外平凡化的丛对构成商 $GL(N,\mathcal X)/GL(N,\mathcal O)$，即仿射 Grassmannian $\mathrm{Gr}_N$；它同构于 $GL(N,\CC((z)))/GL(N,\CC[[z]])$，也同构于 $U(N)$ 的基于点的环路空间。权 $Lw=(m_1,\dots,m_N)$ 对应的修改空间是点 $z^{-m_j}\delta_{ji}$ 在 $GL(N,\mathcal O)$ 下的轨道，这些轨道（Schubert 胞腔）恰好分层整个 $\mathrm{Gr}_N$。}
\QAs{src p.146}{``The latter space is known as the affine Grassmannian $\mathrm{Gr}_N$ for the group $GL(N)$.''}

\QAq{\S10.1 在三维流形上取什么边界条件？}
\QAa{左端固定 $E_-$（即对 $A$ 取 Dirichlet 条件，$\varphi_0$ 自由），右端令 $\varphi_0=0$、$A$ 满足协变 Neumann 条件 $\iota_nF=0$，并且只商去右端非平凡的规范变换。这样定义的椭圆复形的指标给出修改空间的虚维数；无 't~Hooft 算子时该模空间退化为一点。}
\QAs{src p.150}{``we want to allow all possible Hecke modifications $E_+$ that may arise from $E_-$ by solving the Bogomolny equations''}

\QAq{单极子冒泡是什么？}
\QAa{等变瞬子在三维图像中可以把一个量子化单位的电荷集中到一点并使其消失，此即单极子冒泡。原文指出它只能发生在存在 't~Hooft 算子的情形；ADHM 分析表明冒泡把权 $(m+s,m-1)$ 逐步降为 $(m+s-1,m)$，这正是 \S9.2 中 associated 权所描述的退化。}
\QAs{src p.156}{``Monopole bubbling can occur only in the presence of an 't~Hooft operator''}

\QAq{\S10.3 的 Kähler 结构与 moment map 起什么作用？}
\QAa{原文在场空间上引入复结构 $I$（坐标为 $A_z$ 与 $B=A_y-i\varphi_0$）、辛形式与 Kähler 度规；Bogomolny 方程于是写成全纯方程加 moment map 方程 $\mu=0$。$\mu=0$ 还要求边界上 $\Tr\epsilon\varphi_0=0$，这恰好由左端固定规范、右端令 $\varphi_0=0$ 实现，所以原文称这两种边界条件是最简单自然的选择。又因规范变换在左端平凡，$G_\CC$ 自由作用，辛商与复商相同：$Z=Z'$。}
\QAs{src p.161}{``they are the simplest ones that allow us to interpret the Bogomolny equations in terms of a holomorphic equation''}

\QAq{\S10.4 中两个 't~Hooft 算子的 OPE 是什么？}
\QAa{当 $p'\to p$ 时，乘积 $T(Lw,p)\cdot T(Lw',p')$ 展开为 $T(Lw+Lw',p')$ 再加上所有 associated 权 $Lw_i$ 的项。原文说这是 't~Hooft 算符乘积展开的一部分；取 $B,B'$ 为复结构 $I$ 中的复 Lagrangian 子流形（即 $(B,A,A)$ 膜）时，超对称构型空间正是 $Z(LR_1,\dots,LR_k)$，物理态取其上同调。}
\QAs{src p.165}{``this statement is part of the operator product expansion of the operators $T(Lw,p)$ and $T(Lw',p')$''}

\QAq{为什么必须放弃时间反演不变性？}
\QAa{时间反演不变性只适用于所取膜使解满足 $\phi=0$ 的情形；几何 Langlands 最重要的 A-膜（Hitchin 纤维与 c.c. 膜）并不限于 $\phi=0$。因此研究 't~Hooft 算子对它们的作用时必须放弃该假设，改用扩展 Bogomolny 方程。其解沿 $y$ 方向只差规范变换，直到穿过 't~Hooft 算子处发生 Higgs 丛的 Hecke 修改。}
\QAs{src p.179}{``requiring time-reversal invariance puts a severe restriction on the allowed boundary conditions''}

\QAq{Higgs 丛的 Hecke 修改为何是有限的？}
\QAa{全纯丛可以在一点处修改而保持全纯性，但全纯场 $\varphi$ 不能：$\varphi$ 在 $z_0$ 之外给定后，至多只有唯一的全纯延拓。原文算出所需条件数为 $\sum_{i<j}(m_i-m_j)$，恰好等于 (9.25) 给出的修改空间维数，因此 Higgs 丛的 Hecke 修改通常是有限集；当 $\varphi(z_0)$ 正则半单时确实有限，此时 $\varphi$ 在 $Y(Lw,z_0)$ 上定义向量场，其零点就是可能的修改。}
\QAs{src p.180}{``including $\varphi$ generically eliminates some parameters''}

\QAq{典范余各向同性膜是什么？}
\QAa{原文先定义余各向同性子流形：它局部由 Poisson 交换的一阶约束函数定义，维数至少为一半。取支撑为整个 $\MH$、Chan-Paton 为秩一线丛时，A-膜条件 $(\omega^{-1}F)^2=-1$ 给出的复结构可积，且 $\omega-iF$ 必须是全纯辛形式，故这类膜只当 $\omega$ 是全纯辛形式的实部时才可能存在。$\MH$ 上取 $F=(\Im\tau)\omega_J$（此时 $\hat N=I$）所定义的膜称为典范余各向同性膜（c.c. 膜）；它同时是 $(A,B,A)$ 膜，即在复结构 $I$、$K$ 中是 A-膜、在 $J$ 中是 B-膜。}
\QAs{src p.186}{``We will call the A-brane on $M_H$ with $F=(\Im\tau)\omega_J$ the canonical coisotropic brane''}

\QAq{c.c. 膜对应什么 D-模？}
\QAa{c.c. 膜的存在意味着复结构 $K$ 中每个 A-膜都是 $K_M^{1/2}$ 上微分算子 sheaf 的模；任意两个膜的弦构成相应代数上的模，这是开闭 TFT 的公理。对 $\varphi=0$ 的 $(B,A,A)$ 膜，BRST 算符 $Q$ 就是 $\bar\partial$，物理弦 sheaf 是 $K_M^{1/2}$ 的全纯截面 sheaf。type F 膜（Hitchin 纤维加平坦线丛）则给出 $\MH$ 上的 D-模 sheaf。}
\QAs{src p.198}{``it is a magnetic eigenbrane because it is the S-dual of a zerobrane; and it gives rise to a sheaf of D-modules''}

\QAq{一般 $\Psi$ 下还有 c.c. 膜吗？}
\QAa{有。原文证明对任意 $\Psi$ 存在类似物 $B^\Psi_{\mathrm{c.c.}}$，其开弦代数是 $K_M^{1/2}\otimes L^\Psi$ 上的全纯微分算子 sheaf。作为代价，当 $\theta\ne2\pi n$ 时 Hitchin 纤维不再是合法的 A-膜，因为 $B|_Y$ 的周期不是 $2\pi$ 的整数倍；点也永远不可能是 A-膜。}
\QAs{src p.202}{``We will see that an analogue of the c.c. brane exists for all $\Psi$.''}

\QAq{扭化微分算子与普通微分算子如何联系？}
\QAa{对整数 $\Psi=n$，$(B^n_{\mathrm{c.c.}},B_n)$ 的弦态代数恰是 $K_M^{1/2}\otimes L^n$ 上的微分算子；再由扰动论证 $f(\Psi)$ 是线性函数，可推出任意 $\Psi$ 时代数为 $L^{\Psi-h}$ 上的全纯微分算子（$h$ 为对偶 Coxeter 数）。取 $\Psi=h$ 即得作用在函数上的普通微分算子，于是 S-对偶加上 $T^h$ 把零膜映为普通的 D-模。}
\QAs{src p.206}{``if we let $\Psi = h$, the algebra of c.c. strings is the sheaf of ordinary differential operators''}

\QAq{四类膜如何从规范理论得到？}
\QAa{原文列出四类：零膜为 $(B,B,B)$，其镜像是 Hitchin 纤维、为 $(B,A,A)$，c.c. 膜为 $(A,B,A)$，其镜像亦然。四类膜都保两个拓扑超荷（连同其伴随超荷共四个超对称，在规范理论中称 $1/4$ BPS 膜，在 sigma-模型语言中称 $1/2$ BPS 膜）。原文强调这些重要膜都可以在四维规范理论中定义为边界条件，尽管提升并不唯一。}
\QAs{src p.207}{``We will focus on these four kinds of brane and describe the corresponding boundary conditions in the gauge theory.''}

\QAq{零膜 $(B,B,B)$ 的边界条件如何确定？}
\QAa{取 $T=\Gamma^{04}$、$T\lambda|=\lambda|$，得到 $F_{0I}+F_{4I}=0$；再令 $A_0+A_4=0$、其余 $A_I$ 取 Dirichlet 条件，并要求存在超对称态，就要求 Hitchin 方程成立。边界值于是在 $\MH$ 中定义一点 $p$，这正是支撑在 $p$ 的零膜；对不可约 Higgs 丛，这一边界条件还被唯一确定。}
\QAs{src p.213}{``The boundary conditions we have obtained so far correspond to a zerobrane on $M_H$ supported at the point $p$.''}

\QAq{$(B,A,A)$ 的边界条件给出什么？}
\QAa{对 $(B,B,B)$ 作 S-对偶，取 $T'=\Gamma^1\Gamma^2\Gamma^3\Gamma^4$、$T'\lambda|=-\lambda|$：$\varphi_z$ 的特征多项式固定，而决定 $E$ 全纯结构的 $A_z$ 可以涨落。再要求 $\Gamma^{2367}\epsilon=\epsilon$，即得 $D_z\varphi_z=0$（复 Hitchin 方程），膜支撑在 Hitchin 纤维上，是典型的 $(B,A,A)$ 膜。}
\QAs{src p.216}{``the field $A_z$ which determines the holomorphic structure of $E$ is allowed to fluctuate''}

\QAq{$(A,B,A)$ 的边界条件给出什么？}
\QAa{取 $T'=-\Gamma^{15}\Gamma^{2-7}\Gamma^{3+6}$ 并令 $T'\lambda|=\lambda|$，得 $A_z+i\varphi_z$ 固定、$A_{\bar z}$ 自由，从而在 $J$ 中定义复的、中间维数的、$\Omega_J$ 迷向的子流形 $W'$。$W'$ 正是固定全纯结构的平坦 $G_\CC$ 丛，故其上的膜是 $(A,B,A)$ 膜；S-对偶后得到 target 空间为整个 $\MH$ 的 space-filling $(A,B,A)$ 膜。}
\QAs{src p.218}{``So the brane supported on $W'$, with a trivial Chan-Paton bundle, is an (A, B, A)-brane, as expected.''}

\QAq{$(A,B,A)$ 的规范场是否就是 c.c. 膜的数据？}
\QAa{是。对 S-对偶得到的 space-filling 膜，把低能边界条件与混合 Dirichlet--Neumann 条件比较，可定出靶空间规范场 $f=\omega_J$。原文指出这正是 c.c. 膜的靶空间规范场，因此该边界条件所构造的膜就是典范余各向同性膜。}
\QAs{src p.220}{``This is the target-space gauge field for the canonical coisotropic (A, B, A)-brane''}

